% 本文件是待应用的方法模板，不是已执行实验的陈述。
% 应用前核验实际目标、模型绑定、场景数量、协议版本和导出状态。
\subsubsection{L3--来源内容正确性与关键行为保留}
\label{sec:exp-l3-v2}
L3从两个方面评价生成的Detail-UC. 第一部分依据原始来源审计全部最终断言的内容正确性, 第二部分检查预先选定的有限场景中必要行为是否得到保留. L1仅核验断言自身的引用, 而L3允许独立查找原始来源, 因此无引用或引用不当不直接等同于本层内容错误. 本实验不构建全文需求全集, 也不使用原相对完整率RCR.

\noindent\textbf{来源内容审计}\
评价保留各方法最终输出中的原始断言及其顺序, 不在评价阶段拆分或改写. 用例名称和主要参与者属于外部输入, 不计入生成断言分母. 两个Agent分别在匿名条件下核验原始来源, 判断断言是否受到支持、属于目标范围、保持必要的角色与条件关系, 且不存在不必要的重复. 复合断言只有全部实质性内容成立时才判为可保留. 原文查找未获得足够依据的情况记为未决, 不直接判为事实错误.

对每份用例, 设最终断言数为$N$, 已确认可保留数为$K$, 已确认应删除数为$D$, 语义未决数为$U$, 技术未完成数为$F$, 满足$N=K+D+U+F$. 在$N>0$且$F=0$时, 报告已确认可保留比例
\begin{equation}
\mathrm{CP}_{\mathrm{confirmed}}=K/N.
\end{equation}
同时报告未决界限$[K/N,(K+U)/N]$. 该界限不是统计置信区间, 其补数也不等同于错误率. 对空输出或技术未完成的用例不填造比例, 并保留相应数量和状态. 各方法的宏平均仅在同一文档内七种方法共同非空且无技术失败的目标集合上计算, 明确报告该集合大小及其与全部任务的差异.

\noindent\textbf{有限场景检查}\
在评分之前, 根据来源为每个主实验目标选择一个核心成功场景; 来源存在明确影响执行或结果的条件分支时, 再选择一个分支场景. 场景选择者不接触方法输出. 每个场景只登记定义该场景所需的条件、关键动作、结果及其关系, 经独立来源核验后冻结, 供全部方法共用. 场景材料不作为生成方法的额外输入.

评价允许不同措辞及多个槽位共同表达必要行为. 各检查点分别标记为正确保留、部分保留、缺失、关系失真或未决; 程序按预设规则形成场景状态. 核心和分支场景分别报告各状态数量及未建立场景的目标数. 该结果仅衡量所选场景中的行为保留, 不解释为全文需求覆盖率或全局最小充分性.

\noindent\textbf{复核与统计}\
两个主评Agent使用独立会话, 第三个Agent复核分歧并抽检一致项. 复核保留盲判、首评和裁决记录, 不以多数一致替代原文依据. 模型、供应方参数、输入与场景版本以及随机呈现规则均在评分前冻结. 小试只用于检查流程、固定样例判断和费用, 不依据AUCMR的相对得分选择指标或场景.

结果按来源文档分别报告, 以相同目标和场景形成配对比较, 不将断言、评分项或Agent调用次数当成独立项目样本. 本实验的主分析为描述性统计, 不据两份文档推断跨项目显著优势. 未决界限可能改变比较方向时, 保留不确定结论. 任何旧RCR结果均不沿用至本评价.

% 执行器验收通过后生成以下文件；当前没有真实结果，不应在原稿直接启用。
% \input{table_l3_content.tex}
% \input{table_l3_behavior.tex}
% \input{table_l3_errors.tex}
